\documentclass[10pt,a4paper]{amsart}
\usepackage{amsmath,amssymb,amsthm,graphicx}
\usepackage{verbatim}
\newenvironment{reference}{\comment}{\endcomment}
%\newenvironment{reference}{}{}
\newenvironment{slogan}{\comment}{\endcomment}
\newenvironment{history}{\comment}{\endcomment}

% For commutative diagrams we use Xy-pic
\usepackage[all]{xy}

% We use 2cell for 2-commutative diagrams.
\xyoption{2cell}
% Theorem environments.
%
\theoremstyle{plain}
\newtheorem{theorem}[subsection]{Theorem}
\newtheorem{proposition}[subsection]{Proposition}
\newtheorem{lemma}[subsection]{Lemma}

\theoremstyle{definition}
\newtheorem{definition}[subsection]{Definition}
\newtheorem{example}[subsection]{Example}
\newtheorem{exercise}[subsection]{Exercise}
\newtheorem{situation}[subsection]{Situation}

\theoremstyle{remark}
\newtheorem{remark}[subsection]{Remark}
\newtheorem{remarks}[subsection]{Remarks}

\numberwithin{equation}{subsection}

% Macros
%
\def\lim{\mathop{\mathrm{lim}}\nolimits}
\def\colim{\mathop{\mathrm{colim}}\nolimits}
\def\Spec{\mathop{\mathrm{Spec}}}
\def\Hom{\mathop{\mathrm{Hom}}\nolimits}
\def\Ext{\mathop{\mathrm{Ext}}\nolimits}
\def\SheafHom{\mathop{\mathcal{H}\!\mathit{om}}\nolimits}
\def\SheafExt{\mathop{\mathcal{E}\!\mathit{xt}}\nolimits}
\def\Sch{\mathit{Sch}}
\def\Mor{\mathop{\mathrm{Mor}}\nolimits}
\def\Ob{\mathop{\mathrm{Ob}}\nolimits}
\def\Sh{\mathop{\mathit{Sh}}\nolimits}
\def\NL{\mathop{N\!L}\nolimits}
\def\CH{\mathop{\mathrm{CH}}\nolimits}
\def\proetale{{pro\text{-}\acute{e}tale}}
\def\etale{{\acute{e}tale}}
\def\QCoh{\mathit{QCoh}}
\def\Ker{\mathop{\mathrm{Ker}}}
\def\Im{\mathop{\mathrm{Im}}}
\def\Coker{\mathop{\mathrm{Coker}}}
\def\Coim{\mathop{\mathrm{Coim}}}

% Boxtimes
%
\DeclareMathSymbol{\boxtimes}{\mathbin}{AMSa}{"02}

%
% Macros for moduli stacks/spaces
%
\def\QCohstack{\mathcal{QC}\!\mathit{oh}}
\def\Cohstack{\mathcal{C}\!\mathit{oh}}
\def\Spacesstack{\mathcal{S}\!\mathit{paces}}
\def\Quotfunctor{\mathrm{Quot}}
\def\Hilbfunctor{\mathrm{Hilb}}
\def\Curvesstack{\mathcal{C}\!\mathit{urves}}
\def\Polarizedstack{\mathcal{P}\!\mathit{olarized}}
\def\Complexesstack{\mathcal{C}\!\mathit{omplexes}}
% \Pic is the operator that assigns to X its picard group, usage \Pic(X)
% \Picardstack_{X/B} denotes the Picard stack of X over B
% \Picardfunctor_{X/B} denotes the Picard functor of X over B
\def\Pic{\mathop{\mathrm{Pic}}\nolimits}
\def\Picardstack{\mathcal{P}\!\mathit{ic}}
\def\Picardfunctor{\mathrm{Pic}}
\def\Deformationcategory{\mathcal{D}\!\mathit{ef}}
\usepackage{cleveref}
\crefname{originallemma}{original lemma}{original lemmas}
\Crefname{originallemma}{Original lemma}{Original lemmas}
\makeatletter
\newlabel{lemma-ses-artinian}{{0FCH}{source}{}{sourcepointer-1}{}}
\newlabel{lemma-ses-artinian@cref}{{[originallemma][0][]0FCH}{[1][0][]source}{}{}{}}
\makeatother
\crefname{originallemma}{original lemma}{original lemmas}
\Crefname{originallemma}{Original lemma}{Original lemmas}
\makeatletter
\newlabel{lemma-ses-noetherian}{{0FCI}{source}{}{sourcepointer-2}{}}
\newlabel{lemma-ses-noetherian@cref}{{[originallemma][0][]0FCI}{[1][0][]source}{}{}{}}
\makeatother
\title{Objects of finite length and filtrations with simple successive quotients}
\author{The Stacks Project Authors}
\date{stacks-project:homology.tex:local-snapshot-20261012; GNU FDL 1.2 or later}
\begin{document}
\maketitle
\noindent\textit{Editorial note.} All delivered mathematical text is the author's own, copied from the pinned Stacks Project chapter file homology.tex; no statement or proof step was rewritten, reordered or removed.\par\medskip
\noindent\textit{Editorial note.} The proof is the author's complete argument in both directions; the base case of the induction is stated and justified in one sentence, the induction step names the two lemmas it uses and the short exact sequence it applies them to, and no step of the proof is left to the reader.\par\medskip
\noindent\textit{Editorial note.} Every labelled reference occurring inside the counted spans is carried in assembly.external\_references as a non-load-bearing navigation pointer (same-chapter ones at their definition site, cross-chapter ones at the invocation inside the counted span), so no load-bearing mathematics is replaced by a stub. No bibliography citation occurs inside the counted spans.\par\medskip
\noindent\textit{Editorial note.} The same-name check for this card scanned every staged Stacks chapter of this repository -- all sources/stacks-*/pinned-original-source.tex files, the card's own chapter excluded -- and the labels found are recorded in dedup.same\_name\_cross\_chapter on the card.\par\medskip
\noindent\textit{Editorial note.} The preamble is the chapter's own native preamble excerpt already used for this source: the comment-type reference, slogan and history environments (native preamble lines 8-13), the Xy-pic setup (native preamble lines 14-18) and the native theorem declarations with the mathematical macros (native preamble lines 154-217, the actual last line of the file).\par\medskip
\noindent\textit{Allowed standard prerequisites (editorial).}\begin{itemize}
\item Artinian and Noetherian objects. Same chapter, the definitions of an Artinian and of a Noetherian object. Supply the two chain conditions that appear in the statement and are used to make the chain of subobjects exist and terminate.
\item Chain conditions and short exact sequences. Same chapter, Lemma lemma-ses-artinian and Lemma lemma-ses-noetherian. Supply that in a short exact sequence the middle object satisfies a chain condition as soon as the two outer objects do, which is the induction step.
\item Simple objects. Same chapter, the definition of a simple object. Supplies the notion used to describe the successive quotients of the filtration.
\end{itemize}\medskip
\ifdefined\hypertarget\hypertarget{sourcepointer-1}{}\fi
\noindent\textit{Source reference.} Original lemma 0FCH: sources/stacks-homology/pinned-original-source.tex, source lines 2207--2207. Navigation pointer invoked at proof line 2255; the counted span contains the pointer only and the cited result is not reproduced inside it, so the pointer is registered as a non-load-bearing navigation stub in the established form for Stacks material. No bibliography citation is involved.\par\medskip
\ifdefined\hypertarget\hypertarget{sourcepointer-2}{}\fi
\noindent\textit{Source reference.} Original lemma 0FCI: sources/stacks-homology/pinned-original-source.tex, source lines 2218--2218. Navigation pointer invoked at proof line 2255; the counted span contains the pointer only and the cited result is not reproduced inside it, so the pointer is registered as a non-load-bearing navigation stub in the established form for Stacks material. No bibliography citation is involved.\par\medskip
\begin{lemma}
\label{lemma-finite-length}
Let $\mathcal{A}$ be an abelian category. Let $A$ be an object
of $\mathcal{A}$. The following are equivalent
\begin{enumerate}
\item $A$ is Artinian and Noetherian, and
\item there exists a filtration
$0 \subset A_1 \subset A_2 \subset \ldots \subset A_n = A$
by subobjects such that $A_i/A_{i - 1}$ is simple for $i = 1, \ldots, n$.
\end{enumerate}
\end{lemma}

\begin{proof}
Assume (1). If $A$ is zero, then (2) holds. If $A$ is not zero, then
there exists a smallest nonzero object $A_1 \subset A$ by the Artinian
property. Of course $A_1$ is simple. If $A_1 = A$, then we are done.
If not, then we can find $A_1 \subset A_2 \subset A$ minimal
with $A_2 \not = A_1$. Then $A_2/A_1$ is simple. Continuing in this way, we
can find a sequence $0 \subset A_1 \subset A_2 \subset \ldots $
of subobjects of $A$ such that $A_i/A_{i - 1}$ is simple. Since $A$
is Noetherian, we conclude that the process stops. Hence (2) follows.

\medskip\noindent
Assume (2). We will prove (1) by induction on $n$. If $n = 1$, then
$A$ is simple and clearly Noetherian and Artinian. If the result holds
for $n - 1$, then we use the short exact sequence
$0 \to A_{n - 1} \to A_n \to A_n/A_{n - 1} \to 0$
and Lemmas \ref{lemma-ses-artinian} and \ref{lemma-ses-noetherian}
to conclude for $n$.
\end{proof}

\end{document}
